% ==========================================
% BAB V KESIMPULAN
% ==========================================
\cleardoublepage
\chapter{KESIMPULAN}
\label{chap:kesimpulan}

Penelitian ini berangkat dari kebutuhan transformasi digital di lingkungan Kementerian Pertahanan Republik Indonesia dalam menghadapi tantangan pertahanan modern yang semakin kompleks, dinamis, dan dipengaruhi oleh pesatnya perkembangan teknologi informasi.
Hasil analisis menunjukkan bahwa implementasi transformasi digital pertahanan memerlukan pendekatan yang terstruktur untuk mengintegrasikan proses bisnis, data, aplikasi, teknologi, tata kelola, dan keamanan siber.
Selain itu, belum tersedianya kerangka \textit{Enterprise Architecture} yang secara khusus dirancang untuk mendukung kebutuhan pertahanan Indonesia menjadi salah satu kendala dalam mewujudkan interoperabilitas sistem, integrasi data, dan pengambilan keputusan yang cepat serta akurat.

Untuk menjawab permasalahan tersebut, penelitian ini mengusulkan pengembangan TOGAF Pertahanan Indonesia, yaitu \textit{framework Enterprise Architecture} yang diadaptasi dari \textit{TOGAF Architecture Development Method} (ADM) dan disesuaikan dengan karakteristik organisasi pertahanan Indonesia.
Pengembangan \textit{framework} dilakukan dengan pendekatan \textit{Design Science Research Methodology} (DSRM) yang memungkinkan perancangan artefak sebagai solusi terhadap permasalahan organisasi.
\textit{Framework} yang diusulkan mencakup arsitektur bisnis, data, aplikasi, teknologi, tata kelola, serta penambahan fase C5ISR \textit{Architecture} (\textit{Command, Control, Communication, Computer, Cyber, Intelligence, Surveillance, and Reconnaissance}) untuk mendukung kebutuhan \textit{smart defense}, interoperabilitas lintas domain, dan operasi pertahanan modern.

Berdasarkan kajian literatur, analisis kebutuhan organisasi, dan rancangan artefak yang telah disusun, TOGAF Pertahanan Indonesia diharapkan mampu menjadi \textit{blueprint} transformasi digital Kementerian Pertahanan Republik Indonesia yang selaras dengan kebijakan Sistem Pemerintahan Berbasis Elektronik (SPBE) Nasional, mendukung integrasi sistem dan data, meningkatkan kesiapan keamanan siber, serta memperkuat kemampuan pertahanan berbasis informasi.
Selain itu, \textit{framework} yang dihasilkan diharapkan dapat menjadi referensi bagi pengembangan arsitektur digital pada organisasi pertahanan lainnya yang memiliki karakteristik serupa.

Artefak yang dihasilkan akan dievaluasi menggunakan metode \textit{Enterprise Architecture Scorecard}\texttrademark{} dan \textit{Expert Judgement} yang melibatkan akademisi, praktisi \textit{Enterprise Architecture}, serta pakar pertahanan dan teknologi informasi.
Hasil evaluasi tersebut diharapkan dapat menunjukkan bahwa TOGAF Pertahanan Indonesia memiliki tingkat kelengkapan, keselarasan, kelayakan implementasi, serta relevansi yang tinggi terhadap kebutuhan transformasi digital pertahanan nasional, sehingga dapat digunakan sebagai landasan strategis dalam mendukung terwujudnya ekosistem pertahanan digital yang terintegrasi, adaptif, aman, dan berkelanjutan.
